\subsection{What finite source operations must supply for exclusion}
\label{sec:pauli-source-selection-boundary}

Finite-dimensional response, unitary covariance and a quadratic occupation
generator do not alone imply the scalar deformation law used in the
stability theorem. A normalized operational interpretation can replace
that law, but requires an exact identification between a field operator
and an observer operation. The following constructions distinguish these
requirements.

\begin{proposition}[Globally covariant finite occupation cutoff]
\label{prop:covariant-cutoff-statistics-boundary}
Let $\mathcal K=\mathbb C^m$ with $m\geq1$, and
\[
 \mathcal F_{\leq2}
   =\mathbb C\oplus\mathcal K\oplus\operatorname{Sym}^2\mathcal K.
\]
This positive Hilbert space has dimension $(m+1)(m+2)/2$; the two-mode
case has dimension six and the four-mode spin edge has dimension fifteen.
Compress the usual bosonic
creation and annihilation operators to this total-occupation cutoff,
writing $C(v)=P_{\leq2}b^\dagger(v)P_{\leq2}$ and $A(v)=C(v)^\dagger$.
Then $C$ is linear in $v$, the normalized empty vacuum is annihilated
by all $A(v)$, and for every unit $v$,
\[
 N_v=C(v)A(v),\qquad [N_v,C(v)]=C(v).
\]
For every genuine unitary $U$ on $\mathcal K$,
\[
 \Gamma_{\leq2}(U)C(v)\Gamma_{\leq2}(U)^\dagger=C(Uv).
\]
For any Hermitian one-particle $h$, its restricted second quantization
$H=d\Gamma(h)$ obeys
$[H,C(v)]=C(hv)$ and is bounded below. Nevertheless,
$C(v)^2\Omega\ne0$ for every nonzero $v$, and no scalar $q$ satisfies
$A(v)C(v)-qC(v)A(v)=I$ even for one fixed unit mode.
\end{proposition}
\begin{proof}
Annihilation lowers total occupation, so $C(v)A(v)$ is the restriction
of the ordinary bosonic number operator
$d\Gamma(|v\rangle\langle v|)$. It preserves total-occupation sectors
and commutes with the cutoff projection. The ordinary one-particle
commutator identities therefore survive compression. At total occupation
two the creation operator vanishes, so both sides of the claimed raising
commutators vanish there as well. The cutoff also commutes with every
$\Gamma(U)$ and with $d\Gamma(h)$, proving full unitary covariance and
the dynamical identity. The finite spectrum gives
\[
 H\geq 2\min\{0,\lambda_{\min}(h)\}\,I.
\]
The normalized one-mode vectors $|v\rangle$ and $|2v\rangle$ have
number eigenvalues one and two, while
$\|C(v)^2\Omega\|^2=2$. On $|v\rangle$ the proposed scalar law
requires $2-q=1$, hence $q=1$. On $|2v\rangle$ it instead requires
$-2q=1$. More generally its exact commutator defect is
\[
 A(v)C(v)-C(v)A(v)=I-(N_v+I)P_2,
\]
where $P_2$ projects onto total occupation two. This state-dependent
boundary term cannot be replaced by a scalar deformation parameter.
\end{proof}

For a supplied hopping matrix with eigenvalues $\pm\kappa$, the
cutoff Hamiltonian has minimum $-2\kappa$ and admits double occupation
of its negative mode. Thus its covariance, exact quadratic phase law
and stability coexist with failure of exclusion. The construction is a
counterexample to this specified finite implication. It is not a model
of the complete observer axioms: in particular a projective one-particle
transport does not automatically lift to a projective transport on all
occupation sectors, since
$\Gamma(e^{i\theta}I)=e^{i\theta N_{\rm total}}$ need not be scalar.
No source production of the cutoff field operators is asserted.

\begin{proposition}[Normalized operations and an exact occupation generator]
\label{prop:contractive-creation-exclusion}
On a finite-dimensional positive Hilbert space, let $C$ be a creation
operator, $A=C^\dagger$, and $N=CA=A^\dagger A$. Assume
\[
 0\leq N\leq I,\qquad [N,C]=C.
\]
Then $C^2=0$. Consequently, for a complex-linear family $v\mapsto C(v)$
satisfying these conditions for every unit $v$, one has
$\{C(v),C(w)\}=0$ for all $v,w$. No scalar deformation law is required.
\end{proposition}
\begin{proof}
Choose an orthonormal eigenbasis of $N$, with eigenvalues
$\lambda_i\in[0,1]$. The commutator identity gives
\[
 (\lambda_i-\lambda_j)C_{ij}=C_{ij}.
\]
Every nonzero matrix entry therefore maps a zero-eigenvalue vector to
a one-eigenvalue vector. A second application vanishes, so $C^2=0$.
Scaling extends this conclusion from unit vectors to all modes, and
expanding $C(v+w)^2=0$ gives the anticommutation identity.
\end{proof}

Under the same phase law the converse holds: if $C^2=0$, then
$CC^\dagger C-C^2C^\dagger=C$ gives $CC^\dagger C=C$.
Thus $C$ is a partial isometry and a contraction. Within this phase-law
class, contraction is equivalent to exclusion. Identifying an unscaled
field operator with a normalized observer operation therefore carries
the full selection content; it is not a harmless change of units.

This conclusion alone does not construct ordinary many-fermion Fock
space. On $\mathbb C\Omega\oplus\mathcal K$, the operators
$C(v)=|v\rangle\langle\Omega|$ satisfy the contraction and phase
hypotheses for every unit $v$, but allow at most one particle in total.
For $\dim\mathcal K\geq2$, their mixed anticommutator is
\[
 \{C(v)^\dagger,C(w)\}
   =\langle v,w\rangle|\Omega\rangle\langle\Omega|
      +|w\rangle\langle v|,
\]
which is not generally $\langle v,w\rangle I$. Additional structure
is required to recover the full canonical field algebra.

For a linear family on $m$ orthonormal modes, the exclusion test reduces
to a finite set of exact operator identities. It suffices to check the
contraction and unit phase conditions on every $e_i$ and every
$(3e_i+4e_j)/5$ with $i<j$. The proposition gives $C(e_i)^2=0$;
expanding the square at the pair vector then gives
$\{C(e_i),C(e_j)\}=0$. Linearity proves exclusion for every mode.
To obtain full mixed CAR by the paired-instrument route, it suffices to
check its completeness identity at the basis vectors, these real pairs,
and the imaginary pairs $(3e_i+4i e_j)/5$. The two pair identities force
the symmetric and antisymmetric cross terms to vanish. These are exact
matrix identities, not conclusions from finitely sampled outcome
frequencies.

Basis-mode checks alone are insufficient. Set $\sigma_+=|1\rangle\langle0|$.
Two commuting hard-core bit raisers
$C_1=\sigma_+\otimes I$, $C_2=I\otimes\sigma_+$ each satisfy
the contraction and phase conditions, but
$C=(3C_1+4C_2)/5$ has
$C^2|00\rangle=(24/25)|11\rangle\ne0$ and fails the unit phase law.
Inserting the fermionic sign on the second raiser changes no
occupation-basis transition probability. Coherent superposition tests,
or source-derived identities that imply them, must therefore carry the
relative-sign information.

The bound $N\leq I$ follows if the unscaled removal operator $A$ is
itself a Kraus operator of a trace-nonincreasing observer operation.
Equivalently, $A$ is a contraction; its adjoint $C$ is a contraction too.
This is an operational sufficient condition for exclusion only after the
same operator's quadratic form $A^\dagger A$ is identified with the
exact unit-step occupation generator. A normalized probability effect
and a particle-number operator are not interchangeable merely because
both are positive.

A stronger operational interface supplies the mixed relation as well.
If a complete two-outcome instrument has exactly the adjoint-paired Kraus
operators $A(v),C(v)$ for each unit $v$, trace preservation gives
\[
 C(v)A(v)+A(v)C(v)=I.
\]
With the exact phase law and linearity, the preceding proposition and
complex polarization recover all canonical anticommutation relations.
The adjoint pairing, exhaustive two-outcome list and normalization for
every input state are essential hypotheses; they are not consequences
of probability normalization alone.

The cutoff example demonstrates the distinction. For a unit mode its
operator $N_v+A(v)C(v)$ has eigenvalues at most three. Consequently
\[
 E_{\rm remove}=\tfrac14N_v,\qquad
 E_{\rm add}=\tfrac14A(v)C(v),\qquad
 E_{\rm idle}=I-\tfrac14\bigl(N_v+A(v)C(v)\bigr)
\]
are positive and sum to identity. Kraus operators $A(v)/2$, $C(v)/2$
and $\sqrt{E_{\rm idle}}$ define a normalized instrument which retains
the possibility of double occupation. Its removal effect is $N_v/4$,
not the occupation generator $N_v$. In particular the phase commutator
with that effect has coefficient $1/4$, not one. The extra idle outcome
and the operator normalization cannot be suppressed in transferring an
instrument to the exclusion theorem.

Projective readout has a different algebraic role. A L\"uders outcome
with Kraus operator $P=P^\dagger=P^2$ has
$[PP^\dagger,P]=0$, so a nonzero projector cannot satisfy the unit-step
creation law. For the binary coordinate effects, an alternative pair
$|1\rangle\langle0|$, $|0\rangle\langle1|$ is adjoint-paired,
complete and satisfies the single-mode exclusion laws, but swaps the
outcome-certain states rather than preserving them. The declared phase
instrument and its swap-twisted comparison therefore illustrate
different operations, despite carrying the same outcome effects.
Their matrix-level semantics do not select a native particle-creation
instrument. The distinction is explicit in
\path{Lean/EventAlgebra/LuedersPhaseInstrument.lean} and
\path{Lean/EventAlgebra/SourcePhaseInstrumentOutcomeBridge.lean}.

There is a separate obstruction in the finite history-count interface of
\path{Lean/QFT/SourceHistoryGNSDynamics.lean}. Its eight three-bit histories
carry the repair count
\[
 Q(x_0,x_1,x_2)
  =\mathbf1_{x_0\ne x_1}+\mathbf1_{x_1\ne x_2}.
\]
The eigenspaces of its diagonal operator have dimensions two, four and
two at eigenvalues zero, one and two respectively. No operator $A$ on
this full eight-dimensional carrier can obey both
$Q=A^\dagger A$ and $[Q,A^\dagger]=A^\dagger$. Indeed the adjoint
commutator sends $A$ from the eigenvalue-one space into the eigenvalue-zero
space. On the former space $A^\dagger A=Q$ would make $A$ an isometry,
which cannot map four dimensions into two. This excludes identifying
that particular repair count with one quadratic mode occupation; it
excludes neither multiple mode generators nor different source
observables or larger carriers.

Counted transitions can supply a weaker positive statement. If a record
basis has count labels $n_x$ and an operator $T$ has nonzero entry
$T_{yx}$ only when $n_y=n_x+1$, then
$[\operatorname{diag}(n_x),T]=T$ exactly, for arbitrary transition
amplitudes. This follows entry by entry. It does not identify
$\operatorname{diag}(n_x)$ with $TT^\dagger$: that additional equality
constrains transition magnitudes, ranks and degeneracies. The native
history packet fixes its diagonal counts and empirical state. Its
extension to the complete matrix algebra does not authenticate these
occupation-changing operations or their quantum instrument semantics.

Finite covariant response therefore does not force the scalar deformation
law, while normalized operations
can imply exclusion under an exact quadratic occupation/phase
identification. The stated history-count interface supplies neither that identification
nor an exhaustive adjoint-paired quantum instrument. These finite
results leave open an observer-native construction of the required
operations, a different quantization mechanism, and the physical
spin--statistics attachment.
